\documentclass{article}
\usepackage[utf8]{inputenc}
\usepackage{amsmath}
\usepackage{geometry}
\geometry{margin=2cm}
\begin{document}

\section*{Detecção de Genes Expressos pela RT-PCR}

\subsection*{Interpretação e Estratégia}
Esta questão pede para identificar o tipo de genes detectados pela técnica RT-PCR, que utiliza a transcriptase reversa para gerar DNA a partir de RNA. Entender que a RT-PCR mede expressão gênica ativa é fundamental para resolver a questão.

A estratégia consiste em reconhecer:
\begin{itemize}
  \item Que a RT-PCR converte RNA mensageiro em DNA complementar (cDNA).
  \item Que o RNA mensageiro é produzido por genes ativos, em expressão.
  \item Assim, a técnica detecta quais genes estão sendo expressos, não apenas presentes.
\end{itemize}

\subsection*{Conteúdo e Mini Revisão}

\begin{tabular}{|p{4cm}|p{6cm}|p{6cm}|}
\hline
\textbf{Conceito} & \textbf{Ideia-Chave} & \textbf{Aplicação} \\
\hline
RT-PCR & Transforma RNA em DNA complementar para análise & Detecta genes ativos, produzindo RNA mensageiro \\
\hline
Expressão gênica & Genes ativos transcrevem RNA mensageiro & RT-PCR identifica genes ativos pela presença de mRNA \\
\hline
Transcriptase reversa & Enzima que sintetiza DNA a partir de RNA & Fundamental para conversão do RNA em DNA na RT-PCR \\
\hline
\end{tabular}

\bigskip

A técnica RT-PCR inicia pela síntese do DNA complementar a partir de RNA mensageiro, permitindo analisar quais genes estão ativos (expressos). Genes inativos não produzem mRNA e, portanto, não são detectados.

\subsection*{Resolução e Pulo do Gato}

\begin{enumerate}
  \item A RT-PCR começa convertendo o RNA mensageiro em DNA complementar via transcriptase reversa.
  \item O RNA mensageiro está relacionado aos genes ativos, em processo de expressão.
  \item Portanto, o DNA complementar sintetizado indica quais genes estão expressos.
  \item A alternativa correta identifica a técnica como ferramenta para detectar genes em atividade.
\end{enumerate}

\textbf{Pulo do gato:} Lembre que a RT-PCR identifica RNA mensageiro, ou seja, indica a expressão gênica ativa, não apenas a presença do gene.

\subsection*{Armadilhas}

Confundir RT-PCR com PCR tradicional, pensar que RT-PCR detecta todos os genes independentemente da expressão, associar incorretamente a técnica a características genéticas como dominância ou localização cromossômica e limitar seu uso a genes bacterianos ou plasmidiais são os principais equívocos observados.

\subsection*{Código da Questão}
\texttt{CN\_BIOLOGIA\_TECNICA\_001}

\end{document}